\documentclass[11pt]{article}
\usepackage[margin=0.8in]{geometry}
\usepackage[T1]{fontenc}
\usepackage{lmodern}
\usepackage{amsmath,amssymb,amsthm,booktabs,array}
\usepackage[colorlinks=true,linkcolor=blue,citecolor=blue,urlcolor=blue]{hyperref}
\setlength{\parindent}{0pt}
\setlength{\parskip}{0.38em}
\setlength{\emergencystretch}{2em}
\Urlmuskip=0mu plus 2mu
\newtheorem{theorem}{Theorem}[section]
\newtheorem{lemma}[theorem]{Lemma}
\newtheorem{corollary}[theorem]{Corollary}
\newcommand{\N}{\mathbb{N}}
\newcommand{\Z}{\mathbb{Z}}
\newcommand{\Q}{\mathbb{Q}}
\newcommand{\R}{\mathbb{R}}

\title{Batch 15: Motzkin convolutions and bounded zero-sum arrays}
\author{BlakeStone\\\small Proofs developed with ChatGPT}
\date{November 28, 2026}
\begin{document}
\maketitle
\begin{abstract}
We prove the recorded A045994 recurrence from its literal Motzkin
convolution and derive the compatible generating function, correcting a
different expression printed in the source. We then count the literal
bounded zero-sum arrays underlying A202253 and A202254, proving their
floor and parity formulas, recurrences, and source-indexed generating
functions. These are two mathematical families covering three OEIS
entries. Complete human proofs accompany ten checked Lean modules.
Classical ingredients and inspected prior work are credited; the bounded
searches do not establish worldwide priority.
\end{abstract}

\section{A045994: a recurrence for a Motzkin convolution}

Let $m_0=m_1=1$ and
\[
 m_{r+2}=m_{r+1}+\sum_{k=0}^{r}m_km_{r-k}\qquad(r\ge0)
\]
be the Motzkin numbers. The defining recurrence of A045994 is
\begin{equation}\label{eq:a045994-source}
 a_0=1,\qquad
 a_n=m_n+\sum_{k=1}^{n-1}m_ka_{n-k-1}\quad(n\ge1).
\end{equation}
In particular, $a_0,\ldots,a_5=1,1,3,7,18,47$.
We prove the recurrence conjectured by R.~J.~Mathar in 2013
\cite{a045994-source}.

\begin{theorem}\label{thm:a045994}
For every integer $n\ge5$, the sequence defined by
\eqref{eq:a045994-source} satisfies
\[
\begin{aligned}
0={}&(n+1)a_n+2(-2n-1)a_{n-1}+2(-n+3)a_{n-2}\\
 &+(11n-19)a_{n-3}+(11n-27)a_{n-4}+3(n-3)a_{n-5}.
\end{aligned}
\]
\end{theorem}

\begin{proof}
Work in the formal power-series ring $\mathbb Q[[x]]$, and put
$M=\sum_{n\ge0}m_nx^n$ and $A=\sum_{n\ge0}a_nx^n$.
The two defining recurrences give
\begin{equation}\label{eq:a045994-gfs}
 M=1+xM+x^2M^2,\qquad A=M+x(M-1)A.
\end{equation}
The second equality follows directly from the Cauchy product:
the coefficient of $x^n$ in $x(M-1)A$ is
$\sum_{k=1}^{n-1}m_ka_{n-k-1}$.
It gives $(1+xA)M=(1+x)A$, or equivalently
$A=M/(1+x-xM)$; the denominator has constant term one.

Define
\[
 p=x^4+3x^3+2x^2-x,\qquad q=x^2+2x-1.
\]
Eliminating $M$ from \eqref{eq:a045994-gfs} yields
\begin{equation}\label{eq:a045994-quadratic}
 pA^2+qA+1=0.
\end{equation}
For an explicit elimination certificate, set $U=1+xA$ and
$V=(1+x)A$. Polynomial expansion gives
\[
\begin{aligned}
pA^2+qA+1={}&
 U^2\{x^2M^2+(x-1)M+1\}\\
 &-(UM-V)\{x^2(UM+V)+(x-1)U\}.
\end{aligned}
\]
Both terms on the right vanish by \eqref{eq:a045994-gfs}.

Differentiating \eqref{eq:a045994-quadratic} gives
\begin{equation}\label{eq:a045994-derivative}
 (2pA+q)A'+p'A^2+q'A=0.
\end{equation}
Now put
\[
\begin{aligned}
u&=x-4x^2-2x^3+11x^4+11x^5+3x^6,\\
v&=1-6x+2x^2+14x^3+17x^4+6x^5,\\
w&=1-4x+x^2,\qquad
L=1-6x+8x^2+2x^3-x^4.
\end{aligned}
\]
A second polynomial certificate is
\[
\begin{aligned}
(2pA+q)(uA'+vA-w)
={}&u\{(2pA+q)A'+p'A^2+q'A\}\\
&+L(pA^2+qA+1).
\end{aligned}
\]
The right side is zero by \eqref{eq:a045994-quadratic} and
\eqref{eq:a045994-derivative}. The constant coefficient of
$2pA+q$ is $-1$, so it is a unit. Cancellation proves the exact
differential equation
\begin{equation}\label{eq:a045994-ode}
 uA'+vA=w.
\end{equation}

For $n\ge6$, coefficient extraction in
\eqref{eq:a045994-ode} gives the asserted recurrence: if
$u=\sum_{j=0}^{5}\alpha_jx^{j+1}$ and
$v=\sum_{j=0}^{5}\gamma_jx^j$, its left coefficient is
$\sum_{j=0}^{5}\{\alpha_j(n-j)+\gamma_j\}a_{n-j}$.
The coefficient of $x^n$ on the right is zero.
The remaining endpoint $n=5$ follows from the six initial values:
\[
6\cdot47-22\cdot18-4\cdot7+36\cdot3+28\cdot1+6\cdot1=0.
\]
This proves every index $n\ge5$.
\end{proof}

\paragraph{Correction to the printed generating function.}
The source also prints $1/(1-x(1+x)M)$, which is incompatible with
its defining recurrence. Since $M=1+x+2x^2+O(x^3)$, that reciprocal
has coefficient $8$ at $x^3$, whereas \eqref{eq:a045994-source}
gives $a_3=7$. The proof above therefore uses the literal source
definition and derives its transform. The discrepancy and the
recurrence are one target family.

\paragraph{Prior work and formal scope.}
The Motzkin generating-function equation is classical
\cite{a045994-motzkin}. The source links Kruchinin and Kruchinin's
general composita methods \cite{a045994-kruchinin}; these are
credited as background. Niu's 2026 treatment of A176677
\cite{a045994-niu} is a related application of algebraic elimination
to a different recurrence. No novelty is claimed for that method.
The bounded source and repository search, including inspected
implementation hits, is recorded in \texttt{source\_checks.md};
it is not a proof of global priority.

The Lean development starts from the two literal natural-number
recurrences, proves \eqref{eq:a045994-gfs}, checks both polynomial
certificates, and extracts coefficients for the full range.
Its final theorems are
\texttt{A045994.source\_recurrence} and
\texttt{A045994.printed\_generating\_function\_false}.
No assertion about the separate path interpretation in the source
comments is included.


\section{A202253 and A202254: bounded zero-sum arrays}

For $r\in\{3,4\}$ and $n\ge0$, let $A_r(n)$ denote the number of
ordered integer arrays $(a_1,\ldots,a_r)$ such that
\[
 -n\le a_i\le n,\qquad \sum_{i=1}^r a_i=0,\qquad
 -n\le a_i-a_{i+1}\le n\quad(1\le i<r).
\]
Thus adjacency is along the linear array, without a wraparound edge.
The source sequences A202253 and A202254 begin at $n=1$
\cite{a202253-source,a202254-source}; at $n=0$ their natural
combinatorial extensions both equal one.

\begin{theorem}\label{thm:bounded-zero-sum}
For every integer $n\ge0$,
\[
 A_3(n)=1+2\left\lfloor\frac{2n^2+2n}{3}\right\rfloor,
\]
and
\[
 A_4(n)=
 \begin{cases}
 (22n^3+33n^2+26n+12)/12,&n\text{ even},\\
 (22n^3+33n^2+26n+3)/12,&n\text{ odd}.
 \end{cases}
\]
These prove the recorded formula conjectures of Kimberling and Barker
for the literal array definitions.
\end{theorem}

\begin{proof}
First consider three entries. Substitute $z=-x-y$ in an array
$(x,y,z)$. For $y\ge0$, the entry bounds and the two adjacent-difference
bounds are equivalent to
\[
 y-n\le x\le n-2y.
\]
The interval is nonempty exactly when $3y\le2n$. Put
$m=\lfloor2n/3\rfloor$. The slice at $y=0$ has $2n+1$ points, and
negation bijects every positive-middle slice with its negative-middle
counterpart. Therefore
\begin{equation}\label{eq:three-slices}
 A_3(n)=2n+1+2\sum_{y=1}^{m}(2n-3y+1)
       =2n+1+m(4n-3m-1).
\end{equation}
For $n=3q,3q+1,3q+2$, this gives respectively
\[
 12q^2+4q+1,\qquad12q^2+12q+3,\qquad12q^2+20q+9.
\]
These are exactly the three residue-class evaluations of the asserted
floor expression.

For four entries, write an array as $(a,b,c,d)$ and put $s=b+c$.
Negation pairs the slices at $s$ and $-s$, so first suppose $s\ge0$.
The two outer differences satisfy
$c-d=(a-b)+2s$; since both lie in $[-n,n]$, we have $0\le s\le n$.
Set
\[
 m=n-s,\qquad h=\lfloor m/2\rfloor,\qquad
 \epsilon=m\bmod2,\qquad v=a+c.
\]
Then
\[
 c=s-b,\qquad a=v-s+b,\qquad d=-v-b.
\]
Substituting these expressions into all the entry and difference
bounds gives the following equivalent conditions:
\begin{equation}\label{eq:four-rectangle}
 -h\le b\le n-h-\epsilon,\qquad
 -m\le v\le m,\qquad -m\le b+v\le n.
\end{equation}
The first two intervals contain
$(n+1-\epsilon)(2m+1)$ integer points. To count the upper forbidden
corner, write
\[
 b=n-h-\epsilon-i,\qquad v=m-j.
\]
Since $m=2h+\epsilon$, the inequality $b+v>n$ becomes $i+j<h$.
For $i,j\ge0$ this is a triangle with $h(h+1)/2$ points. All those
points lie in the rectangle: $m\le n$ ensures the horizontal bound,
and the vertical bound follows from $i+j<h\le m$. Conversely,
every rectangle point above the diagonal gives such a pair $(i,j)$.
The map
\[
 (b,v)\longmapsto(n-m-b,-v)
\]
preserves the rectangle and interchanges the upper forbidden corner
with the lower forbidden corner $b+v<-m$. The two corners are disjoint.
Consequently the fixed-$s$ slice has
\begin{equation}\label{eq:four-slice-polynomial}
 F_n(m)=(n+1-\epsilon)(2m+1)-h(h+1)
\end{equation}
points. The central slice $s=0$ corresponds to $m=n$, and all others
are paired by negation, giving
\begin{equation}\label{eq:four-slices}
 A_4(n)=F_n(n)+2\sum_{m=0}^{n-1}F_n(m).
\end{equation}
In particular,
\[
 F_n(2q)=-q^2+(4n+3)q+n+1,\qquad
 F_n(2q+1)=-q^2+(4n-1)q+3n.
\]
For completeness, the two prefix sums used to evaluate
\eqref{eq:four-slices} are
\begin{align*}
 3\sum_{m=0}^{2q-1}F_n(m)
   &=-2q^3+(12n+6)q^2-q,\\
 3\sum_{m=0}^{2q}F_n(m)
   &=-2q^3+(12n+3)q^2+(12n+8)q+3n+3.
\end{align*}
Each follows by summing the displayed quadratic slice polynomials,
or by induction on $q$. Substitution of $n=2q$ and $n=2q+1$ in
\eqref{eq:four-slices} yields the asserted parity formula.
\end{proof}

\begin{theorem}\label{thm:bounded-zero-sum-gf}
As identities of formal power series,
\begin{align*}
 \sum_{n\ge1}A_3(n)x^n
   &=\frac{x(3+3x+2x^2-x^3+x^4)}{(1-x)^3(1+x+x^2)},\\
 \sum_{n\ge1}A_4(n)x^n
   &=\frac{x(7+10x+2x^2+4x^3-x^4)}{(1-x)^4(1+x)}.
\end{align*}
For every $n\ge6$ in the source's indexing,
\begin{align*}
 A_3(n)&=2A_3(n-1)-A_3(n-2)+A_3(n-3)
                    -2A_3(n-4)+A_3(n-5),\\
 A_4(n)&=3A_4(n-1)-2A_4(n-2)-2A_4(n-3)
                    +3A_4(n-4)-A_4(n-5).
\end{align*}
The same recurrences hold at $n=5$ with the natural extension
$A_3(0)=A_4(0)=1$.
\end{theorem}

\begin{proof}
The formulas of Theorem~\ref{thm:bounded-zero-sum} give
\[
 A_3(n+3)-A_3(n)=8n+16,\qquad
 A_4(n+2)-A_4(n)=11n^2+33n+30.
\]
Taking two further finite differences in the first identity and three
in the second proves the stated recurrences. The corresponding
expanded denominator polynomials are
\begin{align*}
 D_3(x)&=1-2x+x^2-x^3+2x^4-x^5
          =(1-x)^3(1+x+x^2),\\
 D_4(x)&=1-3x+2x^2+2x^3-3x^4+x^5
          =(1-x)^4(1+x).
\end{align*}
Write $S_r(x)=\sum_{n\ge1}A_r(n)x^n$, whose constant coefficient is
zero. The recurrences imply $[x^n](D_rS_r)=0$ for all $n\ge6$.
The initial positive-index terms are
\[
 (A_3(1),\ldots,A_3(5))=(3,9,17,27,41),\qquad
 (A_4(1),\ldots,A_4(5))=(7,31,81,171,309).
\]
They give
\[
 D_3S_3=3x+3x^2+2x^3-x^4+x^5,\qquad
 D_4S_4=7x+10x^2+2x^3+4x^4-x^5.
\]
Both denominators have constant coefficient one, so they are
invertible as formal power series; division yields the claims.
\end{proof}

\paragraph{Formal scope and prior work.}
The Lean development defines each sequence as the cardinality of the
literal finite set of source arrays. It proves the signed-slice
bijections, both triangular corner counts, the disjoint partition,
the polynomial sums, and coefficientwise generating-function
identities. The formal generating functions use coefficient zero
equal to zero and source counts at every positive index. All seven
modules are default build roots under pinned Lean~4.19.0 and mathlib.
The floor sequence underlying the three-entry answer is already
A143978 \cite{a202253-floor}; its elementary algebra is prior
background. The existing open-math-agents A202253 file is an admitted
statement, and Sequencelib supplies a synthetic formula definition;
neither inspected file proves the source array count
\cite{a202253-oma,a202253-sequencelib}. This package treats two entries
as one family, with recurrences and generating functions as
consequences of the count formulas. A202253's separate empirical
interpretation by bounded solutions of $w+2x+3y=0$ is outside the
selected claims. The bounded searches recorded with the source
snapshots do not establish worldwide priority.

\section{Verification and scope of the release}

The accompanying repository packages contain all Lean sources, pinned
toolchain and dependency manifests, frozen OEIS source snapshots with
hashes, supplementary exact-arithmetic scripts, build logs, and separate
agent review reports. Each package's complete set of local proof modules
is listed as an explicit default build target. The command
\texttt{lake build} checks the package under Lean~4.19.0 and mathlib
\texttt{c44e0c8ee63ca166450922a373c7409c5d26b00b}.

For this release, each of the two packages passed an author build and
a separate fresh source-only build: four successful package builds in
total. The builds used separate local directories and copied no compiled
proof objects from one another. They shared the same pinned compiler and
external dependency caches. The word ``independent'' describes this
separate compilation and a review by another proof agent, not verification
by a second theorem prover or external peer review.

The printed principal theorem audits use only Lean's standard
\texttt{propext}, \texttt{Classical.choice}, and \texttt{Quot.sound}
axioms. Source checks exclude local admissions, custom axiom declarations,
and \texttt{native\_decide}. Input hashes in the build logs bind the
checks to the published sources. The supplementary finite calculations
help detect source, index, and sign mistakes; the universal mathematical
assertions are established by the human and Lean proofs, with the
formal boundaries stated in the relevant sections.

The counts in this collection refer to mathematical families and selected
recorded assertions. Supporting identities, boundary corrections, and
corollaries are not separately counted as solved open problems. Each
package's \texttt{source\_checks.md} records a bounded prior-work inquiry
and credits known ingredients. Neither the recorded conjecture label nor
those searches establishes current worldwide open status or priority.

\begin{thebibliography}{99}
\raggedright
\bibitem{a045994-source}
The OEIS Foundation, \emph{A045994}; recurrence conjecture by
R.~J.~Mathar, September 27, 2013.
Immutable source commit
\texttt{72c53dd50967ac7166f7d31bf79be95174ec5bc6}.
\url{https://oeis.org/A045994}.

\bibitem{a045994-motzkin}
The OEIS Foundation, \emph{A001006: Motzkin numbers}.
Immutable source commit
\texttt{8496aacd31f383b1dc5d340495c56c84e7514e4b}.
\url{https://oeis.org/A001006}.

\bibitem{a045994-kruchinin}
V.~V.~Kruchinin and D.~V.~Kruchinin,
\emph{Composita and its properties}, arXiv:1103.2582v2, 2013.
\url{https://arxiv.org/abs/1103.2582v2}.

\bibitem{a045994-niu}
T.~Niu, \emph{A short proof of Mathar's 2016 recurrence conjecture
for OEIS A176677}, arXiv:2605.04369v1, 2026.
\url{https://arxiv.org/abs/2605.04369v1}.


\bibitem{a202253-source}
R. H. Hardin; formula conjecture by C. Kimberling; generating function
by C. Barker. \emph{OEIS A202253}, revision 19, March 3, 2018.
\url{https://oeis.org/A202253}.

\bibitem{a202254-source}
R. H. Hardin; formula and generating-function conjectures by C. Barker.
\emph{OEIS A202254}, revision 11, March 3, 2018.
\url{https://oeis.org/A202254}.

\bibitem{a202253-floor}
The OEIS Foundation, \emph{A143978: $a(n)=\lfloor2n(n+1)/3\rfloor$},
revision 56, August 12, 2025. \url{https://oeis.org/A143978}.

\bibitem{a202253-oma}
GuillaumeBAUD-coder/open-math-agents, \emph{Statement of OEIS A202253};
AI formalization attributed in the file to Claude Code and jleguina.
Inspected at commit \texttt{ea1e36f03c18db17e965436584565254807e6c7e};
the theorem is admitted.
\url{https://github.com/GuillaumeBAUD-coder/open-math-agents/blob/ea1e36f03c18db17e965436584565254807e6c7e/statements/oeis-a202253.lean}.

\bibitem{a202253-sequencelib}
W. Moreira and J. Stubbs, with original sequence-synthesis work credited
to T. Gauthier and J. Urban, \emph{Sequencelib synthetic A202253 definition}.
Inspected at commit \texttt{dfb06c18bba196b5aa5bcb3725fe9d201fb9a61a}.
\url{https://github.com/provables/sequencelib/blob/dfb06c18bba196b5aa5bcb3725fe9d201fb9a61a/Sequencelib/Synthetic/A202/A202253.lean}.

\end{thebibliography}
\end{document}
